% Generated draft: requires mathematical and source review.
\documentclass[11pt]{article}
\usepackage[margin=25mm]{geometry}
\usepackage{fontspec}
\setmainfont{Latin Modern Roman}
\usepackage{amsmath,amssymb,mathtools,unicode-math}
\setmathfont{Latin Modern Math}
\usepackage{enumitem,needspace,xurl,hyperref}
\setlength{\parindent}{0pt}
\setlength{\parskip}{5pt}
\setlength{\emergencystretch}{3em}
\newcommand{\N}{\mathbb{N}}
\newcommand{\Z}{\mathbb{Z}}
\newcommand{\Q}{\mathbb{Q}}
\newcommand{\R}{\mathbb{R}}
\newcommand{\C}{\mathbb{C}}
\newcommand{\F}{\mathbb{F}}
\newcommand{\A}{\mathbb{A}}
\newcommand{\PP}{\mathbb{P}}
\newcommand{\E}{\mathbb{E}}
\newcommand{\eps}{\varepsilon}
\newcommand{\dd}{\,\mathrm{d}}
\begin{document}
\section*{UnsolvedMath ID: 30005144}
\textbf{Source problem code:} catalogue entry
\subsection{Definitions and mathematical statement}
Let \((M^3,g)\) be a complete, asymptotically flat, three-dimensional Riemannian manifold without boundary. Assume that its scalar curvature satisfies \(\operatorname{Scal}_g\geq 0\) everywhere and that \((M,g)\) is not isometric to Euclidean \((\R^3,g_{\mathrm{Eucl}})\). For each real number \(c>0\), the supplied record considers the functional
\[
\mathcal A_c(\Omega)=\operatorname{Area}_g(\partial\Omega)-c\operatorname{Vol}_g(\Omega),
\]
where \(\Omega\subset M\) is an admissible region. The record does not specify the class of admissible regions or the interpretation of boundary area for nonsmooth regions. For a smooth region with finite volume and finite boundary area, these terms mean its Riemannian volume and the Riemannian area of its boundary.

Write \(\omega_c\) for the one-parameter min--max value of \(\mathcal A_c\) referred to in the supplied statement. The proposed assertion is that, for every manifold satisfying the stated assumptions and every \(c>0\),
\[
\omega_c<\frac{16\pi}{3c^2}.
\]
This assertion remains conditional on the intended definition of \(\omega_c\): the supplied record gives no sweepout class, parameter-space endpoint conditions, topology of continuity, or admissibility and finiteness conditions defining that min--max value. It also gives no precise definition or decay requirements for asymptotic flatness. These choices must be recovered from the source before the assertion is fully specified.

The retained original passage conjectures that an unidentified condition ``(1)'' holds whenever \(M\) has nonnegative scalar curvature and is not isometric to Euclidean \(\R^3\). The displayed strict inequality above is the supplied registry's interpretation of that condition; the equation and its surrounding definitions are absent from this record. The supplied background describes a conjecture about a suitable minimal hypersurface, whereas the title concerns constant mean curvature and the displayed functional has a positive volume coefficient. That terminology is preserved as an unresolved source qualification, without identifying these formulations.
\subsection{Short English statement}
The supplied conjecture asks whether every complete, boundaryless, asymptotically flat Riemannian three-manifold with nonnegative scalar curvature that is not Euclidean has \(\omega_c<16\pi/(3c^2)\) for every \(c>0\), where \(\omega_c\) is the intended one-parameter min--max value of area minus \(c\) times enclosed volume. The precise min--max construction and asymptotic assumptions are missing, so this is a statement-only draft requiring mathematical review.
\subsection{Sources}
\begin{itemize}
\item Local UnsolvedMath registry: \texttt{data/unsolvedmath/problems.json}, permanent ID \(30005144\), source problem code \texttt{catalogue entry}, title \emph{CMC Min--Max Width Under Nonnegative Scalar Curvature}. The supplied record belongs to the Oberwolfach Reports Open Problems collection and gives proposed year \(2022\); no individual proposer is identified.
\item Supplied source citation: \emph{Geometrie (2023)}, Oberwolfach Reports, DOI 10.4171/owr/2022/28, pp. 1551--1601. \url{https://doi.org/10.4171/owr/2022/28}. The supplied background describes Oberwolfach Reports as documenting mathematical workshops and their open-problem discussions.
\item Supplied literature source, labelled ``OWR report'': \url{https://ems.press/content/serial-article-files/46965}.
\item Supplied literature source, labelled ``Bamler--Li--Mantoulidis, PSC decomposition'': \url{https://doi.org/10.1016/j.aim.2023.109231}.
\item Supplied literature assessment dated \texttt{2026-08-22}: the record labels the problem open, describes the report as recording a conjecture about existence of a suitable minimal hypersurface under nonnegative scalar curvature, and states that cited min--max and stable-minimal-surface results give important special cases but that no general proof of the stated condition was found. This is supplied dated curation, not an independently verified literature assessment; the status and cited claims require independent verification.
\item The supplied registry labels the statement \texttt{corrected\_verified} and reports that the official Oberwolfach report was checked, the question isolated from surrounding exposition, and notation repaired for self-contained publication. These are supplied provenance claims, not verification performed for this draft.
\end{itemize}
\end{document}
